\documentclass[a4j]{jsarticle}
\usepackage[top=20truemm,bottom=20truemm]{geometry}
\usepackage[dvipdfmx]{graphicx}

\begin{document}
\title{\LaTeX{}その２}
\author{20yy31000  愛知　太郎}
\maketitle

\begin{figure}[htb]
  \centering
  \includegraphics[width=3truecm]{./figs/right-hand.eps}
  \caption{フレミング右手の法則}
  \label{fig:fleming-rh}
\end{figure}
\begin{figure}[htb]
  \centering
  \includegraphics[width=3truecm]{./figs/left-hand.eps}
  \caption{フレミングの左手の法則}
  \label{fig:fleming-lh}
\end{figure}
図~\ref{fig:fleming-rh}はフレミングの右手の法則,
図~\ref{fig:fleming-lh}はフレミングの左手の法則%
~\cite{kishimoto_elec+mag06}を表す．

\begin{thebibliography}{9}
\bibitem{kishimoto_elec+mag06}
  岸本忠史，電磁気学入門，培風館，2006．
\bibitem{okumua_latex20}
  奥村晴彦，黒木裕介，LaTeX2e美文書作成入門，改訂第８版，技術評論社，2020．
\end{thebibliography}
\begin{figure}[htb]
  \centering
  \includegraphics[width=0.8\linewidth]{./figs/latex-flow.eps}
  \caption{\LaTeX　の処理手順書}
  \label{fig:latex-flow}
\end{figure}

\end{document}
